\documentclass[11pt,a4paper]{article}
\usepackage[T1]{fontenc}
\usepackage{lmodern,amsmath,amssymb,booktabs,microtype}
\usepackage[margin=25mm]{geometry}
\usepackage[colorlinks=true,urlcolor=blue,linkcolor=blue]{hyperref}
\newcommand{\Tr}{\operatorname{Tr}}
\newcommand{\ad}{\operatorname{ad}}
\title{From de Sitter horizons to bulk particles:\\a Susskind-line research programme}
\author{Alan Day \and Chengheng\thanks{AI research assistant (OpenAI).}}
\date{Research programme and first operator construction --- 5 November 2026}
\begin{document}\maketitle
\begin{abstract}
We organize our de Sitter project around static-patch horizon holography
and infinite-temperature double-scaled SYK. The central objective is to
derive a propagating bulk sector from microscopic horizon dynamics.
Our first research target is the correlation and spectral matrix of
contracted singlet operators in complex SYK, followed by a comparison
with the proposed two-dimensional gauge-theory bulk. We identify the
primary literature, incorporate the corrected location of the entropy
layer, and formulate the finite-system operator problem that will anchor
the subsequent large-$N$ analysis.
\end{abstract}

\section{The physical question we will pursue}
A static observer has access to a region bounded by a cosmological
horizon. The programme is to specify a microscopic Hilbert space,
state, Hamiltonian and observable map that account for this region.
The horizon should encode both its entropy and the particles that can
propagate inside it. We therefore choose the following chain of questions:
\[
 \boxed{\text{horizon degrees of freedom}}
 \ \longrightarrow\ \boxed{\text{singlet dynamics}}
 \ \longrightarrow\ \boxed{\text{bulk fields and interactions}}.
\]
The starting microscopic state is the normalized trace of complex SYK,
$\rho_\infty=I/2^N$, with entropy $S_\infty=N\log2$. The double-scaled
limit takes $N,p\to\infty$ at fixed $\lambda=p^2/N$. Here $p$ counts
creation operators, and also annihilation operators, in the complex
Hamiltonian used below; the total monomial degree is $2p$.

The physical hypothesis to develop is that particular contracted
composites carry the propagating bulk sector. This makes a useful
research objective: obtain their spectral data and reconstruct the
corresponding bulk dynamics. The 2026 proposal \cite{mss26} identifies
singlet spectroscopy as a quantitative test of its gauge-theory
interpretation. We select that test as the first main project.

Our earlier observer and local-quantization calculations remain
supporting material. They supply techniques and comparison models;
the present programme determines which techniques are needed next.

\newpage
\section{Reading the development as a connected argument}
The following reading order follows the physical construction. Exact
passages already inspected and those still to be worked through are
listed in the accompanying literature ledger.

\paragraph{1. The observer's horizon as the holographic screen.}
Dyson--Lindesay--Susskind \cite{dls} formulate static-patch
complementarity. Read this as the motivation for a finite-entropy
observer description. It supplies the question that the microscopic
model must answer.

\paragraph{2. Why infinite-temperature SYK enters.}
Lin--Susskind \cite{ls} introduce tomperature to connect a maximally
mixed microscopic state with a finite physical temperature. Susskind's
scale analysis \cite{scales} and Rahman's dS-JT geometry \cite{rahman}
then provide the microscopic and bulk sides of the proposed dictionary.
Our first reproduction will tabulate their time and energy conventions
and trace the entropy and correlation-time matching.

\paragraph{3. Which operators describe bulk matter.}
The confinement analysis \cite{confine} separates elementary or generic
cord excitations from special contracted composites. The 2024 temperature
paper \cite{temperatures}, especially \S4.2, points to derivative
bilinears as a useful tower. This motivates our operator basis rather
than an average over arbitrary matter cords.

\paragraph{4. The flat-space dynamics.}
Sekino--Susskind \cite{qcd} compare the large-$N$ expansions of SYK and
gauge theory. Miyashita--Sekino--Susskind \cite{mss25} develop the
candidate $1+1$ dimensional gauge-theory interpretation in the Rindler
limit. We will follow the equations from the complex Hamiltonian to
the singlet sector and then to the proposed meson spectrum.

\paragraph{5. Updated scales and a quantitative target.}
The November 2025 correction \cite{correction} moves the entropy layer
from the string scale to a parametrically smaller proper distance.
The October 2026 version of \cite{mss26}, \S4.3, gives the singlet
tower and calls for quantitative spectral agreement. These two papers
set our current rather than historical target.

\paragraph{6. Methods we can reuse.}
Gross--Rosenhaus \cite{gr} extract singlet data and cubic couplings
from SYK correlators. Their calculation offers a useful methodological
template: two-point data organize the particles, higher-point data
organize interactions. We will inspect its kernel and normalization
before building the high-temperature counterpart. Its infrared setting
will be recorded in the comparison appendix.

\newpage
\section{A three-scale picture for the construction}
Use $L$ for the cosmic length and $\ell_s$ for the string length.
The parameter relations organizing the Susskind line are
\begin{equation}
 \frac{L}{\ell_s}\sim p,\qquad
 \frac{L}{\ell_P}\sim N,\qquad
 \lambda=\frac{p^2}{N},\qquad g_{\rm open}^2=\lambda.
\end{equation}
Here $\ell_P$ is the Planck scale in the dS$_3$/reduced-JT dictionary;
the relations are parametric and are not a four-dimensional area law.
The corrected phase-boundary relations \cite{correction} are
\begin{equation}
 T_c\sim \sqrt N\,\ell_s^{-1}
       \sim\sqrt\lambda\,\ell_P^{-1},\qquad
 \rho_{\rm sh}\sim \frac{\ell_P}{\sqrt\lambda}.
\end{equation}
Combining these existing relations gives the working hierarchy
\begin{equation}
 \frac{\ell_s}{\rho_{\rm sh}}\sim\sqrt N,\qquad
 \rho_{\rm sh}\ll\ell_s\ll L
 \quad(N\to\infty,\ 0<\lambda\text{ fixed}).\label{hierarchy}
\end{equation}
The entropy layer and the string-scale region therefore play different
roles in bulk reconstruction. We will track both radial locations when
defining horizon operators and bulk extrapolates.

For orientation, the static-patch metric gives the exact proper-distance
and redshift relations
\begin{align}
 f(r)&=1-r^2/L^2, & \rho(r)&=L\arccos(r/L),\nonumber\\
 \sqrt{f(r)}&=\sin(\rho/L), &
 T_{\rm loc}(\rho)&=\frac{1}{2\pi L\sin(\rho/L)}.
\end{align}
Thus $T_{\rm loc}\sim(2\pi\rho)^{-1}$ near the horizon. The blueshift
at the string scale is of order $p$, whereas at the corrected entropy
layer it is of order $N\sqrt\lambda$. These are immediate consequences
of the geometry and the proposed scale dictionary. They explain why
the older string-layer temperature discussion must be read together
with the later correction.

\paragraph{What will determine the physics.}
Entropy fixes the available microscopic state count. The scale hierarchy
sets the regime in which a locally flat bulk can emerge. The singlet
correlators must then determine what propagates there. The next step is
to specify those correlators before attempting to fit a particle spectrum.

\newpage
\section{First construction: an operator spectral problem}
Fix normalized complex fermions $\{c_i,c_j^\dagger\}=\delta_{ij}$ and
the homogeneous number-conserving Hamiltonian
\begin{equation}
 H_J=\sum_{|I|=|J'|=p}J_{I J'}\,c_I^\dagger c_{J'},
 \qquad J_{I J'}=J_{J'I}^*,\qquad Q=\sum_i c_i^\dagger c_i.
\end{equation}
Ordered multi-indices specify the signs, with $c_{J'}$ the adjoint of
the corresponding ordered creation product. At this algebraic step
the coupling variance is left explicit. For numerical work it will be
matched to \cite{mss25}, (2.2)--(2.3), after fixing the CAR convention.
Set $\mathcal L_J O=[H_J,O]$ and use $O(t)=e^{iH_Jt}Oe^{-iH_Jt}$.
The candidate contracted tower is
\begin{equation}
 M_n(J)=\sum_i c_i^\dagger(i\mathcal L_J)^n c_i .
\end{equation}
The same disorder realization must occur in $M_n$ and $H_J$.
For every normally ordered homogeneous monomial, the CAR give
$\sum_i c_i^\dagger[H_J,c_i]=-pH_J$. Consequently
\begin{equation}
 M_0=Q,\qquad M_1=-ipH_J,\qquad [H_J,M_0]=[H_J,M_1]=0.\label{ward}
\end{equation}
This fixes the phase and normalization of the first two tower elements
in our convention. It also identifies the first possible dynamical
operators at $n\ge2$. Use Hermitian components of these elements as
probes $O_a$.

With $d=2^N$ and $(A,B)=d^{-1}\Tr(A^\dagger B)$, the microscopic
target is
\begin{equation}
 C_{ab}(t)=\overline{d^{-1}\Tr[O_a(t)O_b]},\qquad
 \mathsf S_{ab}(\omega)=\frac{2\pi}{d}\,
 \overline{\sum_{r,s}(O_a)_{rs}(O_b)_{sr}
 \delta(\omega-E_s+E_r)}.\label{spectral}
\end{equation}
For any complex vector $z$, $z^\dagger\mathsf S(\omega)z$ is a positive
measure, since each summand is an absolute square. Its moments obey
\begin{equation}
 G_{ab}=\int\frac{d\omega}{2\pi}\mathsf S_{ab}(\omega)=(O_a,O_b),
 \qquad
 K_{ab}=\int\frac{d\omega}{2\pi}\omega^2\mathsf S_{ab}(\omega)
       =(\mathcal L_J O_a,\mathcal L_J O_b),
\end{equation}
with an ensemble average on the right where appropriate. These standard
spectral identities supply exact sum rules for the first calculation.
The generalized eigenvalues of $Kv=\Omega^2Gv$ measure frequency moments.
Bulk masses will instead require a matching to the Rindler or dS bulk
propagator; the microscopic Hamiltonian generates boost time there.

\newpage
\section{A bilocal bridge to the first dynamical candidate}
The first candidate can be written without numerical time derivatives.
Differentiate the conserved $M_1$ using the product rule:
\begin{equation}
 0=\dot M_1=\sum_i\dot c_i^\dagger\dot c_i+M_2,
 \qquad
 A_2=-M_2=\sum_i[H_J,c_i]^\dagger[H_J,c_i]\geq0.\label{a2}
\end{equation}
Thus $A_2$ is Hermitian in every sample. Center it with the same trace
state, $O_2=A_2-d^{-1}\Tr(A_2)I$. Its nonzero-frequency spectral
weight is our first concrete microscopic target. Positivity of $A_2$
is an operator identity; $O_2$ need not be positive after centering.

Rather than compute each derivative composite independently, organize
the whole tower with a two-time operator,
\begin{equation}
 B_J(\tau;t)=\sum_i c_i^\dagger(t)c_i(t+\tau),\qquad
 \left.\partial_\tau^nB_J(\tau;t)\right|_{\tau=0}=M_n(t).
\end{equation}
Define $b_J(\tau)=d^{-1}\Tr B_J(\tau;0)$ and the connected-in-trace
correlator
\begin{equation}
 F_J(\tau,\sigma;t)=d^{-1}\Tr\!\left[B_J(\tau;t)B_J(\sigma;0)\right]
                  -b_J(\tau)b_J(\sigma).
\end{equation}
At finite $N$ the exact relation is
\begin{equation}
 C_{22}(t)=
 \overline{\left.
 \partial_\tau^2\partial_\sigma^2F_J(\tau,\sigma;t)
 \right|_{\tau=\sigma=0}}.\label{bilocal}
\end{equation}
This turns the first composite two-point problem into a contracted
fermion four-point problem with a specified operator order. The
subtraction is made within each sample before averaging. Equation
\eqref{bilocal} is the interface to a real-time bilocal or ladder
calculation, not yet a solution of that calculation. In the large-$N$
limit its equal-time derivatives must be justified by control of the
frequency moments.

\paragraph{How this guides the next calculation.}
We will derive the kernel for $F_J$ in the chosen high-temperature
limit, then take the derivatives in \eqref{bilocal}. A nonzero spectral
component of $O_2$ is the first object to match to a bulk propagator.
Centering removes the identity expectation; other conserved components
are separated using the energy projectors in the appendix. For a
quadratic number-conserving Hamiltonian, $A_2=c^\dagger h^2c$ is itself
conserved. That limiting case provides a useful check that interactions,
rather than the notation of the tower, generate the dynamics.

\newpage
\section{Research sequence and concrete deliverables}
\paragraph{A. Complete the microscopic--geometric dictionary.}
Read \cite{scales,rahman,confine,mss25,correction} at equation level.
Deliver one convention table: real versus complex $N$, interaction
order, CAR, variance, $H_c/H_s$, proper time, entropy, and the two radial
scales in \eqref{hierarchy}. Reproduce the geometric redshift and the
leading correlation times in a shared convention.

\paragraph{B. Derive the singlet correlation kernel.}
Work through the existing bilocal and ladder methods, starting with
\cite{gr} and the real-time analysis in \cite{qcd}. Construct the
contracted channel before disorder averaging, keeping its $J$ dependence.
Deliver the kernel equation for $C_{ab}$, the conserved modes in
\eqref{ward}, and the first nonzero moments for $n=2,3$, starting with
the bilocal derivative in \eqref{bilocal}.

\paragraph{C. Calculate the spectrum in controlled limits.}
First validate the operator construction with finite-$N$ diagonalization
and the spectral sum rules. Start with small positive $\lambda$ after
double scaling, where \cite{mss25}, \S5.3, proposes the meson comparison.
Then organize corrections in $\lambda$: a fixed-order planar kernel
alone does not sum the finite-$\lambda$ expansion. Keep the Hamiltonian
normalization fixed by the dictionary. Deliver the spectral matrix,
its sample variation, and convergence with operator basis size.

\paragraph{D. Reconstruct particles and then interactions.}
Derive the bulk-to-horizon map in Rindler coordinates and compare the
singlet correlators with the candidate gauge-theory propagators. The
desired outputs are invariant mass ratios and spectral weights. Once
that map is fixed, use three- and four-point functions to determine
couplings and test the interaction interpretation of $\lambda$.

\appendix
\section{Technical notes and present status}
This is the first mainline construction and reading plan, not a completed
singlet-spectrum computation. The normal-ordered identities and spectral
sum rules above are standard CAR and linear-algebra consequences.
The 2026 paper supplies the motivation, not the missing quantitative
spectrum. Low-temperature conformal dimensions from \cite{gr} will be
used as method references, not identified with infinite-temperature
boost frequencies.

For finite samples the group action rotates both the fermions and the
coupling tensor. The contracted family is covariant under this joint
action; a fixed generic $H_J$ has exact $U(1)$ charge conservation, not
an assumed exact $U(N)$ symmetry. Conserved contributions to
\eqref{spectral} sit at $\omega=0$. They can be recorded separately
using $\mathcal P_0O=\sum_E P_EOP_E$, with $P_E$ the full degenerate
energy projector. This subtraction changes the measured probe, not the
state or the Hamiltonian.

At exactly infinite temperature the trace of a commutator of two even
operators vanishes. Equation \eqref{spectral} therefore starts with
trace correlation spectra, not an ordinary bosonic retarded response.
The bulk response and the observer map belong to step D.
The earlier spectral-pentagon note is retained as a separate technical
branch; it is not a prerequisite for this programme.

\newpage
\begin{thebibliography}{99}\small
\bibitem{dls} L. Dyson, J. Lindesay and L. Susskind,
\emph{Is There Really a de Sitter/CFT Duality},
\href{https://arxiv.org/abs/hep-th/0202163v4}{hep-th/0202163v4}.
\bibitem{ls} H. Lin and L. Susskind, \emph{Infinite Temperature's Not
So Hot}, \href{https://arxiv.org/abs/2206.01083}{2206.01083}.
\bibitem{scales} L. Susskind, \emph{De Sitter Space, Double-Scaled
SYK, and the Separation of Scales in the Semiclassical Limit},
\href{https://arxiv.org/abs/2209.09999v1}{2209.09999v1}.
\bibitem{rahman} A. A. Rahman, \emph{dS JT Gravity and Double-Scaled
SYK}, \href{https://arxiv.org/abs/2209.09997v2}{2209.09997v2}.
\bibitem{confine} L. Susskind, \emph{De Sitter Space has no Chords.
Almost Everything is Confined},
\href{https://arxiv.org/abs/2303.00792v1}{2303.00792v1}.
\bibitem{temperatures} A. Rahman and L. Susskind, \emph{Infinite
Temperature is Not So Infinite: The Many Temperatures of de Sitter Space},
\href{https://arxiv.org/abs/2401.08555v3}{2401.08555v3}.
\bibitem{qcd} Y. Sekino and L. Susskind, \emph{Double-Scaled SYK,
QCD, and the Flat Space Limit of de Sitter Space},
\href{https://arxiv.org/abs/2501.09423v2}{2501.09423v2}.
\bibitem{mss25} S. Miyashita, Y. Sekino and L. Susskind,
\emph{DSSYK at Infinite Temperature: The Flat-Space Limit and the
't Hooft Model}, \href{https://arxiv.org/abs/2506.18054v3}{2506.18054v3}.
\bibitem{correction} L. Susskind, \emph{Where is the entropy in
DSSYK--de Sitter? Correction to a Wrong Claim},
\href{https://arxiv.org/abs/2511.10907v1}{2511.10907v1}.
\bibitem{mss26} S. Miyashita, Y. Sekino and L. Susskind,
\emph{Holograms and Standard Models},
\href{https://arxiv.org/abs/2607.05678v2}{2607.05678v2}.
\bibitem{gr} D. J. Gross and V. Rosenhaus, \emph{The Bulk Dual
of SYK: Cubic Couplings},
\href{https://arxiv.org/abs/1702.08016v2}{1702.08016v2}.
\end{thebibliography}
\end{document}
